\begin{table}[htbp]
\centering
\small
\caption{University-Level PCSI by Institutional Characteristics}
\label{tab:inst-means}
\begin{threeparttable}
\begin{tabular}{lcccc}
\toprule
Group & $N$ & Mean & SD & Range \\
\midrule
\multicolumn{5}{l}{\textit{Carnegie classification and control}} \\
\quad Private R1 & 38 & 0.456 & 0.046 & [0.337, 0.534] \\
\quad Public R1 & 86 & 0.427 & 0.051 & [0.328, 0.585] \\
\quad Private R2 & 8 & 0.406 & 0.053 & [0.342, 0.459] \\
\quad Public R2 & 18 & 0.441 & 0.058 & [0.376, 0.580] \\
\multicolumn{5}{l}{\textit{Mission and program}} \\
\quad Land-grant & 47 & 0.420 & 0.047 & [0.338, 0.585] \\
\quad Non-land-grant & 103 & 0.442 & 0.053 & [0.328, 0.580] \\
\quad Medical school & 118 & 0.438 & 0.053 & [0.328, 0.585] \\
\quad No medical school & 32 & 0.422 & 0.047 & [0.342, 0.519] \\
\multicolumn{5}{l}{\textit{Urbanicity}} \\
\quad Rural & 15 & 0.425 & 0.041 & [0.372, 0.519] \\
\quad Semi-urban & 48 & 0.422 & 0.047 & [0.338, 0.539] \\
\quad City & 82 & 0.443 & 0.057 & [0.328, 0.585] \\
\bottomrule
\end{tabular}
\begin{tablenotes}[flushleft]\footnotesize
\item \textit{Notes:} Unit of analysis is the university; PCSI is averaged across each university's policy-in-force years before group statistics are computed. Five universities lack urbanicity data. Rank-based tests are reported in Appendix~\ref{app:inst}.
\end{tablenotes}
\end{threeparttable}
\end{table}
